\documentclass[12pt]{article}
\usepackage{array}
\usepackage{amsmath, amsthm, amssymb}
\usepackage[backend=biber]{biblatex}
\usepackage{enumitem}
\usepackage{fancyhdr}
\usepackage{float}
\usepackage{geometry}
\usepackage{graphicx}
\usepackage{hyperref}
\usepackage[utf8]{inputenc}
\usepackage{tikz}

\usetikzlibrary{trees}
\usetikzlibrary{positioning}

\addbibresource{sources.bib}

\graphicspath{{figures/}}

\renewcommand{\thesection}{\Roman{section}}

\geometry{letterpaper, margin=1in}

\pagestyle{fancy}
\fancyhf{} 

\setlength{\parindent}{0pt}
\setlength{\parskip}{1em}

\lhead{MATH 598}
\rhead{Project Proposal}
\cfoot{\thepage}

\title{Project Proposal\\[1ex] Token vs Context in LLM Representations}
\author{Melody Goldanloo\\[2ex] MATH 598: LLMs}
\date{}

\begin{document}

\maketitle
\thispagestyle{fancy}

\section{Introduction}

Transformer models build representations that use information from the current token and prior context via attention. However, it's unclear how much of a given layer's representation is derived from the current token itself versus information aggregated from earlier tokens (i.e., the context).

Understanding this distinction is important for interpretability. If a representations can be recovered from the current token alone, it likely does not rely on contextual reasoning. On the other hand, representations that cannot be recovered without context likely reflect information that requires context, such as higher-level or relational information.

This project aims to quantify this difference and identify which parts of an LLM's internal representation depend on context. Simply put, the question I'm seeking to answer is: how much of an LLM’s layer activations comes from the current token, and how much comes from context?

\section{Proposed Approach}

The core idea is to estimate how much of an intermediate representation can be predicted using only the current token.

Using a pretrained model such as Gemma and a dataset such as The Pile (e.g., a subset like Pile-10k), I will collect, for each token:
\begin{itemize}
    \item the token embedding
    \item intermediate layer activations
    \item sparse autoencoder (SAE) feature activations
\end{itemize}

I will then train a linear model that maps the token embedding directly to the intermediate representation at a given layer. This effectively ``skips'' the attention mechanism and attempts to reconstruct the representation without any contextual information.

The model will be evaluated using cosine similarity, mean squared error (MSE), and $R^2$. High performance indicates that the representation is largely token-local, while low performance suggests dependence on context.

\subsection*{Feature-Level Analysis}

In addition to predicting raw activations, I will predict SAE feature activations. This allows for a more interpretable analysis of which features depend on context.

Specifically, I will examine:
\begin{itemize}
    \item which features are well-predicted from the token alone
    \item which features are poorly predicted and likely require context
\end{itemize}

\section{Hypothesis and Expected Outcomes}

I hypothesize that:
\begin{itemize}
    \item Early layers will be more predictable from the current token
    \item Later layers will rely more heavily on context
    \item Lower-level features (e.g., lexical or local syntactic information) will be more predictable
    \item Higher-level semantic features will be less predictable and require context
\end{itemize}

If the results show that later layers are still highly predictable from the token, this would challenge the assumption that deeper layers primarily encode contextual information.

Conversely, if even early layers are not predictable, this would suggest that contextual mixing happens earlier than expected.

\section{Evaluation and Experimental Design}

The experiment will proceed as follows:
\begin{enumerate}
    \item Run a dataset through the model and collect token embeddings, layer activations, and SAE activations
    \item Train linear models to predict activations from token embeddings
    \item Evaluate prediction accuracy across layers
    \item Analyze distribution of prediction accuracy across SAE features
\end{enumerate}

Results will be visualized across layers to identify trends in context dependence.

\section{Related Work}

The tuned lens approach \cite{belrose2023tunedlens} studies how intermediate representations can be used to predict model outputs. This project is related but reverses the direction by predicting representations from the input token.

Work on sparse autoencoders \cite{bricken2023monosemanticity} shows that model activations can be decomposed into interpretable features. This project builds on that idea by analyzing which of those features depend on context.

Research on induction heads and in-context learning \cite{olsson2022induction} demonstrates that certain model behaviors rely heavily on context, motivating the need to better understand where and how context is used.

Probing methods have also been widely used to analyze representations in neural networks. For example, \cite{hewitt2019structural} shows that linear probes can recover structured linguistic information from model activations. This supports the use of probing as a tool for understanding what information is encoded in a model.

However, probing has known limitations. Prior work \cite{alain2016understanding} highlights that probe performance does not necessarily reflect what information is directly used by the model, and may instead reflect what is linearly recoverable. This motivates careful interpretation of results in this project.

Finally, prior work on transformer structure \cite{elhage2021framework} suggests that many computations in transformers can be understood in terms of relatively simple components, which supports the idea that certain aspects of representations may be predictable from limited information such as the current token.

\section{Challenges and Potential Failure Points}

A key challenge in this project is interpreting probe performance. Poor prediction performance may indicate true context dependence, but it may also result from the limited expressiveness of the probe model. This makes it important to avoid over-justifying and misinterpreting negative results.

A practical consideration is computational cost. Collecting activations and working with SAE features can be expensive on a local machine. To keep the project manageable, I will initially focus on a smaller model, dataset subset, and a limited number of layers or features. This allows for a controlled and interpretable analysis while remaining computationally feasible. However, I can utilize additional compute resources if needed, which would allow me to scale up parts of the analysis.

Finally, analyzing SAE features in a meaningful way may be difficult, as interpreting what individual features represent is itself an open problem.

\section{Expected Contributions}

This project aims to provide:
\begin{itemize}
    \item A quantitative measure of how context dependence changes across layers
    \item A feature-level analysis of which types of information require context
    \item Insight into how representations are structured in transformer models
\end{itemize}

\section{Conclusion}

This project proposes a method for analyzing how much of a language model's internal representations depend on the current token versus prior context. By combining linear probing with sparse autoencoder features, it aims to provide a more interpretable understanding of where and how context is used within the model.

\printbibliography

\end{document}